The base architecture of Cosys-AirSim, inherited from its predecessor Microsoft AirSim, was robust yet fundamentally limited by a strict single-domain paradigm. A simulation session could be configured to host either Unmanned Aerial Vehicles (UAVs) or Unmanned Ground Vehicles (UGVs), but never both simultaneously. This represented a severe bottleneck for modern robotics research, which increasingly relies on Heterogeneous Multi-Agent Systems (MAS) to solve complex real-world problems.\newline\newline
To bridge this critical gap, this project develops a comprehensive multidomain extension. Hence, this extension transitions Cosys-AirSim from a simple vehicle emulator into a holistic testing ground for advanced multi-agent coordination. At its core, the extension introduces a novel architectural component: the \texttt{SimModeWorldMultidomain} class. This specialized C++ module acts as an orchestrator within the Unreal Engine environment, capable of dynamically allocating dedicated communication ports and handling disparate vehicle physics models concurrently.\newline\newline
The impact of this development is relevant,  particularly within the realm of AI-powered robotics. By enabling an environment where aerial and ground vehicles coexist seamlessly, it unlocks entirely new experimental paradigms. For instance, it becomes possible to train reinforcement learning models where a drone acts as an "eye in the sky," streaming fused sensor data to an edge-AI algorithm that actively guides a ground rover through hazardous terrain. \newline